% Part: computability
% Chapter: computability-theory
% Section: s-m-n

\documentclass[../../../include/open-logic-section]{subfiles}

\begin{document}

\olfileid[pt]{cmp}{thy}{smn}
\olsection{O teorema $s$-$m$-$n$}

\begin{explain}
O próximo teorema é conhecido como ``teorema $s$-$m$-$n$'', por uma
razão que ficará clara em breve. A parte difícil é entender exatamente
o que o teorema diz; depois que você compreender o enunciado, ele
parecerá bastante óbvio.
\end{explain}

\begin{thm}
\ollabel{thm:s-m-n}
  Para cada par de números naturais $n$ e~$m$, existe uma função
  recursiva primitiva~$s^m_n$ tal que, para toda sequência
  $e$, $a_0$, \dots, $a_{m-1}$, $y_0$ ,\dots, $y_{n-1}$, temos
  \[
  \cfind{s^m_n(e, a_0, \dots, a_{m-1})}[n](y_0, \dots, y_{n-1}) \simeq
  \cfind{e}[m+n](a_0, \dots, a_{m-1}, y_0, \dots, y_{n-1}).
\]
\end{thm}

\begin{explain}
É útil pensar em $s^m_n$ como algo que atua sobre \emph{programas}.
Isto é, $s^m_n$ recebe um programa~$e$ para uma função $(m+n)$-ária,
além de entradas fixas $a_0$, \dots, $a_{m-1}$; e devolve um programa
$s^m_n(x, a_0, \dots, a_{m-1})$ para a função $n$-ária dos argumentos
restantes. \iftag{TMs}{Se você pensar em $x$ como a descrição de uma
máquina de Turing, então $s^m_n(e, a_0, \dots, a_{m-1})$ é a máquina
de Turing que, nas entradas $y_0$, \dots,~$y_{n-1}$, acrescenta
$a_0$, \dots,~$a_{m-1}$ no início da cadeia de entrada e executa~$e$.
Cada $s^m_n$ é, então, apenas uma função recursiva primitiva que
encontra um código para a máquina de Turing apropriada.}{}
\end{explain}

\end{document}


